\documentclass[10pt,a4paper]{amsart}
\usepackage[margin=19mm]{geometry}
\usepackage{amsmath,amssymb,mathtools}
\usepackage[scr=rsfs,cal=euler]{mathalfa}
\usepackage{xfrac}
\usepackage[unicode,hidelinks,bookmarks=false]{hyperref}
\newcommand{\ie}{i.e.\ }
\newcommand{\cf}{cf.\ }
\newcommand{\resp}{respectively\ }
\newcommand{\Subsection}{\subsection}
\providecommand{\nobreakdash}{\nobreak\hbox{-}}
\setlength{\emergencystretch}{2em}
\multlinegap0pt
\usepackage{amssymb}
\usepackage{mathtools}
\usepackage{xcolor}
\usepackage[shortlabels]{enumitem}
\newtheorem{lemma}{Lemma}
\title{\bf Mathematical Research Meets Large Language Models}
\author{Paata Ivanisvili, Xinyuan Xie}
\date{Source: arXiv:2605.05193v1 (CC BY 4.0)}
\begin{document}
\maketitle

\noindent\emph{Source and licence.} \bf Mathematical Research Meets Large Language Models. Version arXiv:2605.05193v1 (CC BY 4.0), distributed by arXiv under the Creative Commons Attribution 4.0 International licence, http://creativecommons.org/licenses/by/4.0/, as recorded in the arXiv licence metadata for this exact version and shown on the abstract page https://arxiv.org/abs/2605.05193v1. Full licence notice: \texttt{licenses/arxiv-2605.05193-CC-BY-4.0.txt}.\par\smallskip

\noindent\textit{Editorial scope.} Counted unit: the article's lemma lemma-Lf (source0.tex lines 950--955). The article's complete original proof of it is delivered with it (source0.tex lines 956--973, one \texttt{proof} environment of 18 lines). No mathematics is written, repaired or re-derived: the statement and the proof are the article's own lines, in the article's own notation, and no retained line is substituted or re-flowed.  No further source-local context is needed: every cross-reference of every retained line resolves inside the two delivered spans.  Nothing was omitted from any retained span, so the package carries no omission record.  No retained line contains a citation, so the wrapper carries no bibliography and no bibliography entry.  No retained line contains a figure environment, an image-inclusion command or drawing code, so no image is delivered and the package declares no asset dependency.  Editorial formatting only, in addition: an AMS \texttt{amsart} preamble that restates the article's own declarations and macros used by the retained lines, namely the environments \texttt{lemma} on separate counters, together with the standard packages those lines need.  The retained environments print in this document as Lemma 1 in order of appearance, whereas the article prints the same items with the numbering of its own full document.  The complete native source of the article as distributed in the arXiv e-print for this exact version is bundled unchanged as \texttt{sources/199935/source0.tex}.
\begin{lemma}\label{lemma-Lf}
Let $n\in \mathbb{N}$ be even and $n \geq 4$. Let $f$ be the function discussed above. Then
\[
Lf(x) \leq \frac{2}{n-1} f(x), \, \quad \text{for all } x \in \Omega_n.
\]
\end{lemma}

\begin{proof}
Denote by $S_n(x)=\sum_{i=1}^n a_ix_i$. Note that for $x \in \Omega_n$,
\begin{align*}
    \sum_{i<j} S_n^{(i,j)}(x) &= \sum_{i<j} \big(S_n(x) + (a_i-a_j)(x_j-x_i)\bigr) \\
    & = \binom{n}{2} S_n(x) + \frac{1}{2}\sum_{i\neq j}(a_ix_j -a_ix_i-a_jx_j+a_jx_i) \\
    & = \binom{n}{2} S_n(x) + \frac{1}{2}\sum_{i\neq j}a_ix_j - \frac{1}{2}\sum_{i\neq j}a_ix_i- \frac{1}{2}\sum_{i\neq j}a_jx_j+ \frac{1}{2}\sum_{i\neq j}a_jx_i\\
    & =  \binom{n}{2} S_n(x) + \frac{1}{2}\sum_{i=1}^n (a_i(\sum_{j=1}^n x_j)-a_ix_i) - \frac{1}{2}(n-1)S_n(x) \\ &- \frac{1}{2}(n-1)S_n(x) + \frac{1}{2}\sum_{j=1}^n (a_j(\sum_{i=1}^n x_i)-a_jx_j) \\
    & = \binom{n}{2} S_n(x) - \frac{1}{2}S_n(x) - \frac{1}{2}(n-1)S_n(x)- \frac{1}{2}(n-1)S_n(x) -  \frac{1}{2}S_n(x), \quad\text{since}\sum_{i=1}^n x_i=0 \\
    & = (\binom{n}{2} -n )S_n(x).
\end{align*}
Hence, for $x \in \Omega_n$
\begin{align*}
    Lf(x) &= |S_n(x)| - \frac{1}{\binom{n}{2}}\sum_{i<j}|S_n^{(i,j)}(x)| \\
    & \leq |S_n(x)| - \frac{1}{\binom{n}{2}}|\sum_{i<j}S_n^{(i,j)}(x)| \\
    & = |S_n(x)| - \frac{1}{\binom{n}{2}}|(\binom{n}{2}-n)S_n(x)| \\
    & = \frac{2}{n-1}f(x).
\end{align*}
\end{proof}

\end{document}
